% ============================================================
% Spec-per-Stage：面向 AI 编程的交互界面重构——一个注意力经济学视角
% 编译方式：XeLaTeX（必须，需编译两次以生成目录与引用编号）
% 环境：TeX Live / TeXPage / Overleaf 均可
% ============================================================
\documentclass[11pt,a4paper]{ctexart}

\usepackage{amsmath,amssymb}
\usepackage{geometry}
\geometry{margin=1in}
\usepackage{url}
\emergencystretch=2em
\usepackage[colorlinks=true,citecolor=blue,linkcolor=blue]{hyperref}

\title{Spec-per-Stage: Refactoring the Interaction Interface of AI Programming---An Attention-Economics Perspective}
\author{Yang Zikun\\Independent Researcher}
\date{2026年9月}

\begin{document}

\maketitle

\clearpage

\tableofcontents

\clearpage

\markboth{摘要}{摘要}

\begin{abstract}
AI 编程正在成为软件生产的主导方式，但现有讨论大多聚焦于“模型能生成什么”，而较少审视“人如何向模型下达指令”。本文从注意力经济学的视角，提出一种针对 AI 编程交互界面的重构——Spec-per-Stage（SpS，阶段规范范式）：以一套通用规范充当 AI 编程的“SQL 语法”，将一次开发任务表达为若干条规范语句；每条语句由不变的模板骨架 $\tau$ 与开发者填写的变量细节 $\delta$ 构成；语句在提交生成前经过一道轻量校验（如同 SQL 的语法与类型检查），每个阶段（产品方案设计$\rightarrow$技术方案设计$\rightarrow$产品实现）的出口设置人工评审门禁（如同事务提交）。

本文的三重论证是：第一，注意力是 AI 编程中真正稀缺的资源——模型的上下文注意力随无关标记的注入而稀释，人的评审注意力受带宽约束，而自由提示词把两者的注意力大量消耗在“每次重新决定项目结构”这类不变量上；第二，不变量吸收——通用规范将结构性决策从每次会话的随机变量变为确定性常量，模型与人的注意力因此重新集中于真正携带信息的变量细节，正如数据独立性曾把查询者从访问路径中解放出来；第三，门禁有界——阶段门禁使缺陷在引入后的最近边界被捕获，修复动作退化为“重新生成一条语句”，错误传播与修改半径被结构性压低。

SpS 组织 AI 编程的输入端（指令），与作者先前提出的 Feature-per-File（F2F）范式——组织输出端（文件）——互补衔接：实现语句以 F2F 特征文件为输出单元，复杂任务通过语句分解为多个独立文件的增量拼装。本文同时给出规范语句的参考语法（reference grammar）与规范控制台的交互语义（语句、校验、提交、回滚与历史），并以“活动趋势查询”为运行示例展示完整工作流；最后提出生命周期分工模型：结构生成期（0$\to$1）使用 SpS 快速构建结构完整、功能可用的全新项目，局部微调期（1$\to$1.1）依据相位迁移判据回落到传统 AI 编程工具，规范产物作为两相之间的交接资产。
\end{abstract}

\noindent\textbf{关键词：}AI 编程；注意力经济学；通用规范；声明式界面；规范语句；阶段门禁；生命周期分工；规格驱动开发

\begin{center}
\textbf{Abstract}
\end{center}
\noindent AI programming is becoming the dominant mode of software production, yet most discussions focus on what models can generate rather than on how humans instruct them. From an attention-economics perspective, this paper proposes a refactoring of the interaction interface of AI programming---Spec-per-Stage (SpS): a universal specification grammar plays the role that SQL plays for data access; a development task is expressed as a set of specification statements, each composed of an invariant template skeleton $\tau$ and developer-supplied variable details $\delta$; a lightweight validation pass runs before generation (as SQL syntax and type checking does), and a human review gate stands at the exit of every stage---product design, technical design, and implementation (as a transaction commit does).

Three arguments are offered. First, attention is the truly scarce resource in AI programming: the context attention of models dilutes as irrelevant tokens accumulate, the review attention of humans is bandwidth-bound, and free-form prompting spends much of both on re-deciding project structure---an invariant. Second, invariant absorption: universal specifications turn structural decisions from per-session random variables into deterministic constants, so the attention of both model and human re-concentrates on the variable details that actually carry information, just as data independence once freed query writers from access paths. Third, gate boundedness: stage gates capture defects at the nearest boundary after their introduction, reducing repair to regenerating a single statement, which structurally suppresses error propagation and the modification radius.

SpS organizes the input side of AI programming (instructions) and complements the author's earlier Feature-per-File (F2F) paradigm, which organizes the output side (files): implementation statements emit F2F feature files, and complex projects grow by statement decomposition and incremental assembly of independent files.

The paper also gives a reference grammar for specification statements, specifies the interaction semantics of a specification console (statements, validation, commit, rollback, history), and demonstrates the full workflow on a running example.

It closes with a lifecycle division-of-labor model: the construction phase ($0\to1$) uses SpS to build structurally complete, functional new projects; the fine-tuning phase ($1\to1.1$) falls back to conventional AI tools under a phase-transition criterion, with committed specification artifacts serving as the handoff assets between the two phases.

\noindent\textbf{Keywords:} AI programming; attention economics; universal specification; declarative interface; specification statement; stage gate; lifecycle division of labor; spec-driven development

\clearpage

\section{引言}
本章陈述问题域与两条失控证据，给出以注意力为关键变量的核心论点，并列出本文的主要贡献。

\subsection{被忽视的另一半：AI 编程的交互界面}
AI 编程工具的竞争几乎全部集中在生成端：更强的模型、更长的上下文、更自主的智能体\cite{ref1,ref14,ref15}。而交互端——人如何向模型下达指令——仍停留在“自由文本提示 + 补丁式 diff”的原始形态：开发者在对话框里描述意图，模型改写代码，开发者阅读 diff 判断对错。

两条证据表明这个形态正在失控。其一，GitClear 对大量 AI 辅助仓库的度量显示，代码重复率上升、重构减少、代码质量指标下滑\cite{ref2}——无结构的生成在无结构的指令下持续熵增。其二，开发者的实际痛感并不在“生成得不够多”，而在“改不动、看不过来、错了不知道错在哪”：这三者都不是生成能力问题，而是注意力问题——模型的上下文注意力花在了哪里，人的评审注意力落在了哪里。

现有工具对这一问题的回答是“上下文工程”：让模型看到更多。本文主张相反方向的杠杆：不是让注意力看到更多，而是让人与模型把注意力花得更少、更准。稀缺资源的浪费不在模型能力，而在交互界面的组织方式。

\subsection{核心论点：通用规范是 AI 编程的 SQL，注意力是关键变量}
关系模型的历史为本文提供了一个精确的类比。在 SQL 之前，数据访问是过程性的：使用者必须写出导航路径——索引、游标、循环——语言的大部分注意力消耗在“怎么取”上。Codd 的数据独立性\cite{ref3}与 SEQUEL\cite{ref4}把访问路径交给系统，语法固定为少数子句，使用者的全部语言注意力集中于“要什么”。SQL 语法因此成为数据访问的“通用提示词”：它不随任务而变，变化的只是语句中携带的业务细节。

本文的论点是：AI 编程正处于“过程性导航”阶段。开发者用自由文本反复声明项目结构、技术栈、分层方式、风格约定——这些是“怎么建”的访问路径；而真正稀缺的资源——模型的上下文注意力\cite{ref10,ref11}与人的评审注意力——被这些不变量大量占用。与 SQL 同构的解法是：一套通用规范（specification grammar）充当 AI 编程的 SQL 语法，开发者的注意力只花在语句的变量部分——那些真正区分本项目与其他项目的细节上。

注意力是这个范式的关键变数。范式优劣的判据不是生成代码的行数，而是单位有效产出所消耗的模型注意力与人的注意力。自由提示词是一场注意力拍卖：每次会话都重新竞价结构性决策；SpS 的目标是把这些决策从拍卖台上撤下来，变成语法的常量。

\subsection{主要贡献}
\begin{enumerate}
\item 提出并论证“注意力是 AI 编程的稀缺资源”：从模型端（上下文注意力的稀释效应）与人端（评审注意力的带宽约束）给出双重证据，指出自由提示词的浪费主要发生在不变量上；
\item 提出 SpS 范式：通用规范（语法层）、轻量校验（解析层）、阶段门禁（事务层）三个机制，给出规范语句的形式化定义与参考语法（reference grammar）；
\item 给出 SQL 类比的系统映射与规范控制台（specification console）的交互语义——语句、校验、提交、回滚与历史；
\item 论证 SpS 与 F2F\cite{ref20} 的互补性：SpS 组织输入端（指令），F2F 组织输出端（文件）；复杂任务通过语句分解为多个独立特征文件的增量拼装，错误传播与修改半径在两端同时被有界化；
\item 分析范式的成本与任务边界：面向正向构建，不面向探索式原型与深度重构；
\item 提出生命周期分工模型：0$\to$1 结构生成期使用 SpS，1$\to$1.1 局部微调期依相位迁移判据回落到传统 AI 编程工具，规范产物作为交接资产，结构性变更触发重入。
\end{enumerate}

\section{背景}
本章沿四条线索梳理知识基础：声明式接口的历史回报、规格驱动实践的现状与缺口、LLM 注意力的实证特征，以及阶段门禁与缺陷经济学的经典结论。

\subsection{声明式接口：数据独立性的启示}
Codd 的关系模型\cite{ref3}把数据访问分离为逻辑意图与物理路径两层，查询优化器负责后者；SEQUEL\cite{ref4}进一步把意图的表达固定为一种近乎自然语言的语法。Parnas 的信息隐藏判据\cite{ref5}——把易变的决策封装为模块的秘密——是同一思想在模块层的表述。本文把这一判据搬到提示词层：通用规范把“不易变的决策”封装为语法（模板骨架与常量层），把“易变的决策”留作槽位。声明式的回报在两个时代同构：SQL 时代是优化器接管访问路径，SpS 时代是模型接管结构决策——前提是结构决策必须先被语法确定下来。

\subsection{规格驱动开发：先例与缺口}
规格化文档有长期标准传统：IEEE 29148\cite{ref17}规定了需求文档应当写什么。近年的规格驱动工具（GitHub Spec Kit\cite{ref18}、AWS Kiro\cite{ref19}）把规格作为生成的中间产物，验证了“先规格后代码”路线的可行性；反向方向——从自然语言到结构化查询的语义解析\cite{ref16}——则从数据侧印证了“结构化语法优于自由文本”这一判断。但现有实践存在三个缺口：其一，规格未被当作有语法的“语句”——可校验、可分解、可提交、可回滚；其二，注意力未被当作显式的设计变量——模板的取舍缺乏统一的经济学判据；其三，阶段之间缺少事务语义——评审通过与否对流程状态的影响没有被形式化。SpS 对这三个缺口给出统一回答。

\subsection{LLM 注意力的实证特征}
Transformer 的注意力机制\cite{ref10}是有限资源：上下文越长，模型对关键信息的利用越不均匀。“Lost in the Middle” 效应\cite{ref11}表明，无关标记的注入会稀释模型对相关信息的位置召回与任务准确率；提示工程综述\cite{ref12}系统化了“如何写提示”，但将模板视为提示技巧，而非语法。对本文的推论有二：上下文中每一个含混或无关的标记都在消耗预算；结构性决策若以含混方式反复出现，其成本不止于长度，更在于方差——每次会话对同一结构问题给出不同答案，下游产物随之漂移。人的注意力服从同样的经济学：评审者面对无结构的指令史与不限大小的 diff，注意力没有落点，评审质量随体量衰减。

\subsection{阶段门禁与缺陷经济学}
Royce 的阶段化模型\cite{ref6}与 Boehm 的缺陷成本曲线\cite{ref7}给出经典结论：缺陷发现得越晚，修复成本越高。Naur 的“编程即理论构建”\cite{ref8}提醒：程序的本质是开发者头脑中的理论，而非代码文本本身；把人完全移出回路的自主智能体\cite{ref13}恰恰切断了理论的在场。SpS 的立场是折中的：把人放回少数关键节点（阶段门禁），同时用结构化产物与模板化槽位，把“人在场”的注意力成本压到最低——人评审的不是自由文本，而是结构化文档中的变量槽。

\section{自由提示词的失效模式：注意力去哪儿了}
本章把自由提示词工作流的失效归结为三种模式——结构性重复、注意力稀释、错误晚发现——并以规范语句的形式化收束。

\subsection{结构性重复：每个会话都在重新发明语法}
在自由提示词下，一类信息被反复声明：项目结构、技术栈版本、目录布局、命名约定、接口风格、错误处理策略。它们或者每次会话重新描述，或者被复制进越来越长的“规则文件”（如全局系统提示、仓库规则文件）。这是人类时代习惯在提示词层的延续——人类开发者只对代码执行 DRY，对指令从不执行。

F2F\cite{ref20} 曾论证“复用经济学的反转”：当生成成本趋零，成本的主导项从“写出来”转移到“注意力耗掉”。提示词层的结构性重复正是该反转在输入端的表现：重复的不是代码，而是注意力。每次会话对同一批结构问题的重新声明与重新求解，都是对模型注意力与人的阅读注意力的重复征税。

\subsection{注意力稀释：不变量挤占变量}
把一条自由提示词的信息分解为三类：不变量 $I$（结构性决策，跨任务不变）、变量 $V$（任务细节，跨任务变化）、噪声 $N$（含混、冗余与自相矛盾的表述）。现有工作流的典型比例是 $I$ 与 $N$ 挤占 $V$。

模型端的机制是明确的：长而含混的提示迫使模型在生成时同时求解两类问题——“该遵循什么结构”与“该实现什么细节”。前者本是常量，却以自然语言的方差形式出现；上下文中无关与含混标记的积累进一步稀释对关键细节的召回\cite{ref11}。结果是生成分布的质量重心偏移：结构性决策方差大，细节准确率下降。

人端的机制同样明确：评审者面对的是一段对话史与一组 diff，而非结构化产物。他必须先在头脑中重建“模型被要求了什么”，才能判断“模型做了什么”。评审注意力的带宽\cite{ref8,ref9}被重建过程占用，真正用于判断细节的份额缩小。复杂度没有被消除，只是被转移给了注意力最昂贵的环节——人。

\subsection{错误晚发现：跨阶段放大}
软件开发天然分三个阶段：产品方案设计（做什么、给谁做）、技术方案设计（怎么组织、用什么）、产品实现（写成代码）。自由提示词工作流把三个阶段压扁成一场连续对话：意图的偏差、方案的偏差与实现的偏差在代码阶段才一起显形，且 LLM 以近乎相同的成本与置信度实现一个错误的设计\cite{ref14,ref15}——错误一旦进入实现阶段，修复半径由代码组织方式决定（F2F\cite{ref20} 已论证分层组织会放大修改半径）。

自由提示词的问题在更上游：它让错误更晚才被看见。没有阶段产物，就没有阶段边界；没有阶段边界，就没有“在此处捕获”的位置。缺陷成本曲线\cite{ref7}在 AI 时代不仅成立，而且被生成速度放大：人类时代错误实现的代价是工时，AI 时代的代价是“在错误地基上高速叠加更多错误实现”。

\subsection{形式化：规范语句与提示词方差}
\textbf{定义 1（通用规范）} 通用规范是一个三元组 $(T,G,V)$。$T$ 为模板集，每个模板 $\tau$ 定义阶段类型与槽位集合；$G$ 为项目常量层，记录跨语句共享的结构决策（技术栈、架构风格、命名与风格约定、质量约束）；$V$ 为校验规则集，定义语句与产物必须满足的静态约束。

\textbf{定义 2（规范语句）} 一条规范语句是二元组 $s=(\tau,\delta)$，其中 $\tau\in T$ 为模板骨架（不变的语法），$\delta$ 为开发者对 $\tau$ 所定义槽位的赋值（变量的细节）。语句是 AI 编程的最小工作单元。

\textbf{定义 3（结构决策熵）} 对任务类 $C$，设 $S(c)$ 为完成任务 $c$ 所需的结构决策集（架构、组织方式、风格选择）。在自由提示词 $p$ 下，模型对 $S$ 的后验分布随措辞变化，$H(S\mid p)>0$；在模板 $\tau$ 下，$S$ 的取值由 $\tau$ 与 $G$ 完全确定，$H(S\mid \tau)=0$。

\textbf{命题 1（不变量吸收）} 将自由提示词替换为规范语句后，模型与人的注意力预算中分配给结构决策 $S$ 的份额趋于下界，边际注意力重新集中于 $\delta$；生成分布的方差相应地从结构层转移到细节层，即生成质量的不确定性只存在于开发者真正需要决策的地方。

\textbf{依据} $H(S\mid\tau)=0$ 消除了结构层的全部含混；长上下文稀释效应\cite{ref11}意味着消除无关与含混标记对细节召回的改善是超线性的；人端评审对象从“对话史 + diff”变为“槽位赋值表”，重建成本（3.2 节）归零。

\subsection{小结}
自由提示词的三种失效——结构性重复、注意力稀释、错误晚发现——共享同一根源：把不变量当变量处理。第 4 节给出的范式用一个动作回应：把不变量吸收进语法，让注意力只花在变量上。

\section{SpS 范式：把提示词变成语句}
本章给出 SpS 范式的三个机制、规范语句的参考语法、规范控制台的交互语义、与 F2F 的衔接方式，以及 SQL 类比的系统映射。

\subsection{核心思想：三个机制}
SpS 由三个机制构成，分别对应语法的三个层次：

\textbf{机制一：通用规范（语法层）。}一套覆盖三个标准阶段的模板：$\tau_P$（产品方案设计）、$\tau_D$（技术方案设计）、$\tau_I$（产品实现）。模板定义固定骨架与槽位：$\tau_P$ 的槽位如目标用户、问题陈述、范围边界、成功标准、非目标；$\tau_D$ 的槽位如架构风格、数据访问方式、接口契约、异常策略；$\tau_I$ 的槽位如输出单元、依赖语句、验收清单。项目常量层 $G$ 承载跨语句不变量：技术栈与版本、架构风格（如 F2F）、命名约定。模板通过快捷方式实例化——如同在查询工具中键入片段模板，开发者敲入 /pd、/td、/impl 即可生成对应语句的骨架，注意力直接进入填槽。

\textbf{机制二：轻量校验（解析层）。}语句在提交生成前经过一道静态校验：槽位完整性（必填槽位非空）、一致性（$\delta$ 与常量层 $G$ 无冲突）、契约闭合（语句引用的上游产物与接口均存在且可解析）。校验失败即拒绝生成——如同 SQL 解析器在执行前拒绝语法错误。这一机制的成本极低（纯静态检查），却把最廉价的一类失败从“生成后人工发现”前移到“生成前机器发现”。

\textbf{机制三：阶段门禁（事务层）。}每个阶段的出口设置人工评审门禁：产物通过评审，则提交（commit）进入下一阶段；否则回滚（regenerate），重新生成该阶段的语句产物。阶段门禁给予流程事务语义：任何时刻，项目状态都由“已提交的阶段产物”唯一确定，未提交的生成结果不进入项目历史。

三个机制的分工是清晰的：通用规范消除结构方差（命题 1），轻量校验前移廉价失败，阶段门禁有界化错误传播（5.2 节）。

\subsection{规范语句的参考语法}
模板与快捷方式是通用规范在人机界面上的形态；本节给出其背后的抽象语法——规范语句的参考语法（reference grammar）。它不是一门待实现的语言，而是语句形态的非规范性（non-normative）定义：任何工具，只要其指令形态满足该语法，即为 SpS 兼容——正如 SQL 的前端既可以是控制台，也可以是可视化查询构建器，语法的有效性不依赖于具体界面。

\begin{table}[htbp]
\centering
\caption{规范语句的参考语法（EBNF，非规范性）}
\label{tab:grammar}
\begin{verbatim}
<console>    ::= <session> { <statement> }
<session>    ::= ".use" "spec:" <version> "project:" <identifier>
<statement>  ::= <directive> <assignment>+
<directive>  ::= "/pd" | "/td" | "/impl"
<assignment> ::= <slot> ":" <payload>
<slot>       ::= <identifier>
<payload>    ::= <value> { ";" <value> }
\end{verbatim}
\end{table}

其中，$\langle$session$\rangle$ 载入通用规范版本与项目常量层 $G$；$\langle$directive$\rangle$ 对应三类阶段模板 $\tau_P$、$\tau_D$、$\tau_I$；合法 $\langle$slot$\rangle$ 集合由所选模板定义；$\langle$payload$\rangle$ 即开发者填写的变量细节 $\delta$。

语法的设计原则只有一条：自由度被刻意压缩到唯一的位置。directive 与 slot 的组合由模板约束——语法承担不变量；开发者唯一自由书写的区域是 payload——注意力载荷。这与 SQL 严格同构：SELECT$\ldots$FROM$\ldots$WHERE 的子句骨架是常量，变化的只是字面量与绑定变量；骨架的固定不是限制，而是把使用者的全部注意力推向值的手段。4.5 节的映射表将这一同构展开为逐项对照。

轻量校验（机制二）在语法上的落点由此明确：槽位完整性即必填 $\langle$assignment$\rangle$ 非空；一致性即 $\langle$payload$\rangle$ 与常量层 $G$ 无冲突；契约闭合即语句引用的上游产物均处于已提交状态。校验是语法的静态语义，而非附加的工艺。

\subsection{规范控制台：语句、事务与历史}
SpS 的工作界面是规范控制台（specification console），其交互语义直接类比 SQL 控制面板：语句是工作单元，会话历史是语句序列，提交/回滚是事务动作，项目状态是已提交产物的集合。控制台会话的一个片段如下（沿用 F2F\cite{ref20} 的活动趋势查询作为运行示例，通用规范版本 v0.3）：

\begin{table}[htbp]
\centering
\caption{规范控制台会话示例（spec v0.3）}
\label{tab:console}
\small
\begin{tabular}{@{}l@{}}
\texttt{> .use spec:v0.3 project:activity-board} \quad \# 载入通用规范与常量层 $G$ \\[3pt]
\texttt{> /pd} \quad \# 实例化产品方案设计模板 \\
\texttt{PD-001 ($\tau_P$: 产品方案设计)} \\
\hspace*{2em}\texttt{目标用户: 运营团队} \\
\hspace*{2em}\texttt{问题: 需要查看近 30 日用户活动趋势} \\
\hspace*{2em}\texttt{范围: 单看板页; 只读; 无导出} \\
\hspace*{2em}\texttt{成功标准: 打开即见趋势; P95 查询 $<$ 1s} \\
\hspace*{2em}\texttt{非目标: 权限细分; 移动端} \\
\hspace*{2em}\texttt{[validate] 槽位 5/5 完整; 与 $G$ 无冲突 -> pass} \\
\hspace*{2em}\texttt{[gate] 人工评审 -> commit (2026-09-17)} \\[3pt]
\texttt{> /td} \quad \# 技术方案设计, 消费 PD-001 \\
\texttt{TD-001 ($\tau_D$: 技术方案设计)} \\
\hspace*{2em}\texttt{架构风格: F2F 特征文件 [20]} \\
\hspace*{2em}\texttt{数据访问: SQL, 只读从库} \\
\hspace*{2em}\texttt{接口契约: GET /api/trend?range=30d -> TrendResponse} \\
\hspace*{2em}\texttt{[validate] 契约可解析; 上游引用闭合 -> pass} \\
\hspace*{2em}\texttt{[gate] 人工评审 -> commit} \\[3pt]
\texttt{> /impl} \quad \# 产品实现, 消费 TD-001 \\
\texttt{IMPL-001 ($\tau_I$: 产品实现)} \\
\hspace*{2em}\texttt{输出单元: ActivityTrendQuery.java (F2F 单文件)} \\
\hspace*{2em}\texttt{依赖语句: TD-001} \\
\hspace*{2em}\texttt{验收清单: 契约一致; 静态检查通过; 趋势正确性用例} \\
\hspace*{2em}\texttt{[validate] 静态检查通过; 与 TD-001 契约一致 -> pass} \\
\hspace*{2em}\texttt{[gate] 人工评审 -> commit -> 拼装进项目} \\
\end{tabular}
\end{table}

三个要点值得强调。其一，历史即设计日志：语句序列完整记录了“为什么是这个设计”——每条决策都有语句可查，这是对话式工作流无法提供的可审计性。其二，回滚是局部的：评审不通过时，重生成的是一条语句，而不是一场对话的上下文。其三，控制台不隐藏模型：填槽与评审之间，模型仍然负责生成文档与代码——SpS 改变的是指令的形态，不是取消生成。

\subsection{语句分解与拼装：与 F2F 的衔接}
复杂 SQL 任务通过拆解为多条语句完成；同理，复杂项目通过拆解为多条规范语句完成。分解的单位与拼装的单位需要精确对应，否则语句产物会在集成时重新纠缠。这正是 F2F\cite{ref20} 的位置：F2F 以特征为最小自治单元，单文件承载从入口、业务处理到数据访问的完整链路，文件间零耦合，维护采用整体替换。

SpS 与 F2F 的衔接是双向咬合的：

\textbf{输出单元对齐：}实现语句 $\tau_I$ 的输出单元即 F2F 特征文件。一条实现语句生成一个自治文件，文件的可替换性使“评审不通过$\rightarrow$重生成整条语句”成为零副作用的操作——替换式维护在输入端有了对应物：语句重生成。

\textbf{拼装顺序无关：}F2F 的零耦合保证多个已提交的特征文件可以按任意顺序拼装进项目，语句的完成顺序不约束集成顺序；接口契约在校验层闭合（机制二），集成期不需要新的协商。

\textbf{修改半径双端有界：}输入端，缺陷至多存活到下一个阶段门禁（5.2 节）；输出端，缺陷的修复半径至多是一个特征文件。两端共同把“一次错误的代价”从不可预知压到有界。

这一衔接自然引出一个生命周期问题：项目建成之后，SpS 与既有 AI 编程工具如何分工？第 6 节给出形式化的回答——两个相位、一个迁移判据与一套交接协议。

\subsection{SQL 类比的系统映射}
表~\ref{tab:mapping} 给出两个体系的完整映射。

\begin{table}[htbp]
\centering
\caption{SQL 世界与 SpS 世界的系统映射}
\label{tab:mapping}
\small
\begin{tabular}{|l|l|}
\hline
\textbf{SQL 世界} & \textbf{SpS 世界} \\
\hline
SQL 语法 & 通用规范 $(T,G,V)$ \\
\hline
一条查询语句 & 一条规范语句 $s=(\tau,\delta)$ \\
\hline
字面量 / 绑定变量 & 开发者细节 $\delta$（槽位赋值） \\
\hline
数据字典 / schema & 项目常量层 $G$ \\
\hline
语法与类型检查 & 轻量校验（生成前） \\
\hline
查询计划 & 技术方案（$\tau_D$ 产物） \\
\hline
EXPLAIN 后执行 & 门禁评审后生成 \\
\hline
事务 COMMIT / ROLLBACK & 阶段提交 commit / 语句重生成 regenerate \\
\hline
结果集 & 阶段产物（文档 / 特征文件） \\
\hline
复杂任务拆为多条 SQL & 复杂项目分解为多条语句 \\
\hline
多结果集的应用层拼装 & 独立特征文件按契约拼装（F2F 单元） \\
\hline
查询控制台 / workbench & 规范控制台 \\
\hline
方言（MySQL / PostgreSQL） & 规范方言（项目族模板变体） \\
\hline
\end{tabular}
\end{table}

映射在两处刻意中止，这同样有信息量。其一，SQL 有优化器在语法与执行之间自动选择路径，SpS 在技术方案阶段之后设有人工门禁——因为“软件的结构”不像访问路径那样有单一最优解，它携带业务判断，不应被自动化吞掉。其二，SQL 语句之间通常独立，规范语句之间存在消费关系（$\tau_D$ 消费 $\tau_P$ 产物）——这使 SpS 更像带依赖的批处理语言，也正因如此需要门禁而非仅仅需要解析器。

\section{性质与讨论}
本章陈述范式的三条核心性质，划定其与自由提示词工作流的边界，并诚实列出成本与局限。

\subsection{性质一：注意力集中性}
命题 1 在两端各有可操作的含义。模型端：结构决策的方差归零后，上下文中不再出现“自相矛盾的结构指令”，稀释源被移除，细节召回的改善超出模板文本本身的长度收益\cite{ref11}。人端：评审对象从自由文本变为槽位赋值表，评审动作退化为核对——“这些细节是否是我要的”，而非“这份产物整体是否可信”。前者是抽样，后者是通读；注意力消耗相差一个量级。

需要澄清的是：注意力集中性不承诺总注意力下降（模板本身占据固定开销），它承诺的是注意力的方差与错配下降——注意力从“重新决定已知的事”转移到“决定未知的事”。这正是把注意力作为关键变数的含义：范式优化的目标是注意力分配的质量，而非生成产物的数量。

\subsection{性质二：错误早绑定与传播有界}
\textbf{命题 2（传播有界）} 在阶段门禁流水线 $\Pi=(P\to D\to I)$ 中，第 $k$ 阶段引入的缺陷至多存活到第 $k$ 阶段的出口门禁；修复动作是重生成该阶段语句，不触碰已提交的上游产物。若缺陷逃逸到第 $k+1$ 阶段，其修复需要跨阶段回溯，期望成本由修改半径与重审半径之积决定。

Boehm 曲线\cite{ref7}在此获得一个新的放大因子：LLM 实现错误设计与实现正确设计的成本近乎相同\cite{ref14,ref15}，因此“晚发现的错误”不再以工时计价，而以“已高速生成的错误产物体积”计价。门禁的价值不在评审本身，而在给错误一个尽早显形的固定位置。配合输出端的 F2F 单元（4.4 节及\cite{ref20}），即使缺陷逃逸到实现阶段，修复半径仍被压缩到单文件替换。

\subsection{性质三：评审友好性}
\textbf{命题 3（评审有界）} 自由智能体会话的评审对象是跨文件、大小不可预知的 diff；SpS 的评审对象由模板结构保证有界：阶段 $k$ 的评审 = 一份模板化文档（$\tau_P$、$\tau_D$ 产物）或一个特征文件（$\tau_I$ 产物）加一组槽位级差异。

这一性质直接回应 GitClear 度量揭示的失控模式\cite{ref2}：无结构生成 + 弱评审导致重复率上升与熵增。评审友好性使“每阶段人工 review”从理想主义口号变成可持续的日常动作——评审带宽是整个范式的吞吐上限，命题 3 保证该上限不被产物体积击穿。

\subsection{与自由提示词工作流的边界}
SpS 不是自由提示词的替代，而是其上的一个结构化层。以下场景自由提示词仍然占优：探索式原型（结构本身是待发现物，模板会过早锁死搜索空间）；单文件小修（三阶段流水线的固定开销超过收益）；创意性生成（写作、设计草稿，无“结构正确性”可言）。判断准则与命题 1 一致：任务中不变量占比越高、阶段链越长、错误传播代价越大，SpS 的收益越大。全新项目的正向构建——本文的靶心场景——三项指标同时取最大值。需要说明的是，本节判据回答的是任务维度（哪些任务适合 SpS）；项目时间维度的分工——建成之后如何回落、何时重入——由第 6 节的生命周期分工模型回答。

\subsection{成本与局限}
\textbf{模板僵化与规范锁定。}模板在消除结构方差的同时也会消除结构多样性；模板盲区内的好设计会被模板排除。缓解：模板版本化，方言允许项目族变体（表~\ref{tab:mapping}），且锁定只作用于正向构建期——微调期不受约束（第 6 节）。

\textbf{小任务开销。}三阶段流水线对琐碎任务是负收益。缓解：5.4 节的判断准则；控制台允许语句级跳过（单阶段语句）。

\textbf{规范自身的治理。}通用规范是需要演化的人工制品：模板改版、常量层升级、槽位增删都会引入迁移问题，如同 schema 迁移。这是被 SQL 类比照亮的一个成本：语法不是免费的，但语法的收益随语句数量摊薄。

\textbf{门禁的等待成本。}阶段门禁引入人工等待。缓解：门禁异步化（提交后继续草拟下一条语句，评审通过前不生成下游产物），以及评审本身的槽位化（5.1 节）把单次评审成本压低。

\textbf{对生成模型的不可见性。}SpS 在交互界面层运作，不改变模型权重；模板带来的收益依赖模型对模板的遵循度。低能力模型可能连模板一致性都无法维持——但这恰恰说明 SpS 与模型能力是互补而非替代关系：模型越强，语法释放的注意力越多。

\section{生命周期分工模型}
SpS 与传统 AI 编程工具的关系不是替代，而是生命周期上的分工。本章将这一分工形式化：以结构决策熵（定义 3）随时间的演化为参量，项目生命周期划分为两个相位，各有其占优工具；相位之间由一个可操作的迁移判据与一套交接协议连接。

\subsection{两个相位}
\textbf{相位一：结构生成期（0$\to$1）。}项目从零推进到“结构完整、功能可用”。此相位的特征是结构决策熵处于峰值：架构、组织、风格等每一个结构决策都未定，不变量尚未沉淀，任何缺陷的传播半径都由尚未成形的结构决定。这是 SpS 的作用域：通用规范吸收不变量（命题 1），阶段门禁有界化错误传播（命题 2），语句分解配合 F2F 单元完成拼装（4.4 节）。5.4 节的三项判据——不变量占比、阶段链长度、错误传播代价——在此相位同时取最大值。

\textbf{相位二：局部微调期（1$\to$1.1）。}功能可用之后的演化：缺陷修复、局部增强、参数与文案调整。此相位的特征恰好反转：结构已冻结为 F2F 特征文件，单次修改的修改半径有界——至多一个文件，整体替换\cite{ref20}；修改不引入新的结构决策，熵增量近似为零。此时三阶段流水线的固定开销（模板填充、双重校验、门禁等待）超过其方差消除收益，传统 AI 工具（对话式、补丁式）反而占优：它们以文件为上下文单元的会话形态恰好匹配“修改半径 $\le 1$”的任务；F2F 的特征文件为这类工具提供了安全的修改单元——一次微调至多影响一个文件，文件即边界。

\subsection{相位迁移判据}
两个相位之间的迁移由一个可操作的判据给出。

\textbf{判据（相位迁移）。}对下一次修改 $m$，设 $\Delta H(S\mid m)$ 为其引入的新增结构决策熵，$r(m)$ 为其预期修改半径。当 $\Delta H(S\mid m)\approx 0$ 且 $r(m)\le 1$ 个 F2F 单元时，$m$ 属于局部微调，应交由传统工具在文件级上下文中完成；当 $\Delta H(S\mid m)$ 显著为正（新增跨文件特征、架构调整、技术栈变更）时，$m$ 是一次结构性变更，应重入 SpS 会话，重新走语句与评审。

判据把“项目建成后如何继续”从口味问题变成可检验的边界问题，并预示生命周期的循环形态：微调期的持续累积（$1\to1.1\to\ldots$）最终触发结构性变更，重入新一轮语句化构建（$0'\to1'$）。两个相位的系统对比见表~\ref{tab:phases}。

\begin{table}[htbp]
\centering
\caption{两个相位的系统对比}
\label{tab:phases}
\small
\begin{tabular}{|l|l|l|}
\hline
\textbf{维度} & \textbf{相位一：结构生成期（0$\to$1）} & \textbf{相位二：局部微调期（1$\to$1.1）} \\
\hline
结构决策熵 & 峰值，$H(S\mid p)>0$ & 熵增量 $\Delta H(S\mid m)\approx 0$ \\
\hline
修改半径 & 未定，由结构形态决定 & 有界，$r(m)\le 1$ 个 F2F 单元 \\
\hline
占优工具 & SpS（规范控制台） & 传统 AI 工具（对话式、补丁式） \\
\hline
工作单元 & 规范语句 $s=(\tau,\delta)$ & 特征文件（文件即边界） \\
\hline
核心机制 & 语法吸收不变量 + 阶段门禁 & 文件级上下文 + 整体替换 \\
\hline
相位迁移 & commit 阶段产物后进入相位二 & 结构性变更（$\Delta H>0$）时重入相位一 \\
\hline
\end{tabular}
\end{table}

\subsection{交接协议}
分工不是断崖：两相之间有明确的交接资产，SpS 会话提交的产物在微调期继续服役。

\textbf{规范即上下文。}$\tau_D$ 产物（架构风格、接口契约、异常策略）是微调期最精确的 grounding 材料——传统工具的会话以文件为上下文，而规范文档本身就是文件，可直接引用。相比从代码反推设计，规范把“为什么是这样”显式携带进微调期。

\textbf{常量层即约束。}项目常量层 $G$（技术栈、命名、风格约定）在微调期充当轻量护栏：即使是对话式工具，也可将 $G$ 注入会话上下文，抑制微调漂移。

\textbf{历史即审计。}语句序列（设计日志）为微调提供决策考古层：修改一个特征文件前，可查其生成语句与上游契约（4.4 节的依赖关系），避免微调破坏既有契约。

\subsection{与“取代论”的分别}
生命周期分工是本文与既有主张的一个明确分别：SpS 不宣称自己是更好的 Cursor 或 Claude Code，也不宣称微调期应继续留在语句界面里——正相反，它把“何时离开”写进了范式本身。这与 SQL 世界的形态一致：没有人用查询控制台写数据清洗的一次性脚本，也没有人因此否认 SQL。工具的作用域声明是语法成熟的标志，而非弱点。

\section{扩展讨论：开放问题}
本节给出三个超出本文论证范围、但决定范式命运的问题。

\textbf{规范即数据。}语句序列是结构化、可 diff、可统计的记录。这使开发过程本身第一次成为可分析的数据集：结构决策的演化史、缺陷被捕获的门禁位置分布、槽位填充与返工率的相关性——注意力经济学的命题由此从论证变为可度量。一个自然的度量是注意力集中比 $\kappa$：评审与生成中分配给变量 $\delta$ 的注意力份额；SpS 的实证议程即验证 $\kappa$ 与交付质量的正相关。与之配套的是定义 3 的操作化：结构决策熵可经模板覆盖率近似测量——一次会话中由模板与常量层确定的结构决策占比越高，$H(S\mid\tau)\to 0$ 的承诺兑现得越充分；生命周期分工（6.2 节）的 $\Delta H(S\mid m)$ 同样可循此路径估计。

\textbf{方言与生态。}SQL 的方言分裂提示了通用规范的可能命运：不同项目族、不同组织的模板变体终将分化。分化未必是坏事——方言竞争是语法演化的机制；但跨方言的语句迁移成本值得预研。

\textbf{与 F2F 的合流。}SpS 组织输入端、F2F 组织输出端，两者共享同一经济学根基：人类时代的效率产物（自由提示词、水平分层）在生成成本趋零的时代失效。合流后的完整方法论可以概括为一句话：输入声明式，输出自治式——人写的每条指令是一条有语法的语句，机器产出的每个单元是一个有边界的文件。

\section{结论}
本文的工作可以概括为一句话：把 AI 编程的输入端从自由文本重构为有语法的语句，并界定这套语法在项目生命周期中的作用域。

本文从注意力经济学的视角，提出面向 AI 编程交互界面的重构范式 Spec-per-Stage：通用规范充当 SQL 语法，规范语句表达任务，轻量校验前移廉价失败，阶段门禁有界化错误传播；输入端与 F2F 的输出端范式衔接，并以生命周期分工模型界定了其与传统 AI 编程工具的作用域边界，共同构成面向 AI 编程时代的完整方法论。

对 AI 编程的失控症状\cite{ref2}，本文给出的诊断不是模型不够强，而是界面组织不当：自由提示词把稀缺的注意力拍卖给了不变量。修复的动作与关系模型当年相同——声明式化。SQL 把“怎么取”交给优化器，把注意力还给“要什么”；SpS 把“怎么建”交给语法与模型，把注意力还给“建什么”。

AI 编程的下一个杠杆，可能不在更强的模型，而在更好的语法。把不变量吸收进规范，把注意力还给注意力该在的地方。

\clearpage
\bibliographystyle{unsrt}
\bibliography{references-sps}

\end{document}
